\chapter{Preliminary Work: From Toy Models to a Minimal Model}
\label{ch:development}\label{sec:development}

This chapter records how the questions of Chapter~\ref{ch:intro} took
their final form: the initial framing, the exploratory models built in
the first months of the project, the conceptual error those models
exposed, and the corrections from supervision that fixed the object of
study. We record this material because each failed formulation
constrained the next one, and the sequence explains why the final model
looks the way it does. All preliminary code, interactive laboratories,
and notes are archived in a dedicated preliminary-work repository
(\texttt{Score\_Diffusion}, directory \texttt{Toy-models}), kept
separate from the reproducibility package of Appendix~\ref{app:repro}
and available together with it.

\section{The initial framing: dynamic objects}\label{sec:dev-initial}

The project began from a simple observation about the diffusion
literature. Standard diffusion models generate \emph{static} objects; the
data we wanted to understand are \emph{dynamic}: a data point is a
sequence of frames indexed by an internal time $u = 0, 1, \dots, K-1$ (a
video, a rotating shape, a stochastic trajectory). The naive treatment
embeds the whole trajectory in $\R^{dK}$ and learns a single monolithic
score, which ignores the fact that consecutive frames are strongly
correlated. The initial research hypothesis was correspondingly
ambitious: that the joint score field of a dynamic object has an internal
structure of its own, concretely that consecutive frame-blocks of the
score might be related by a fixed map (a ``propagator''),
\begin{equation}
  \score_{:,u}(x, t) \;\approx\; \mathcal{P}_t\big( \score_{:,u-1}(x, t) \big),
  \label{eq:dev-propagator}
\end{equation}
which a learning system could exploit by learning one block and one map
instead of $K$ blocks. Two distinct times appear already in this
formulation, and keeping them apart became a permanent discipline of the
project: the \emph{internal} time $u$ of the sequence, and the
\emph{diffusion} time $t$ of the noising channel, which acts
independently on every frame.

\section{The toy models}\label{sec:dev-toys}

The exploratory phase produced seven small, self-contained interactive
models (browser-based laboratories with exact formulas inside, no neural
networks), built to make score fields visible before attempting theorems.
In the numbering of the archive:

\begin{itemize}
  \item \emph{TM1: the OU channel on one frame.} A single univariate
  distribution transported by the forward OU flow, with the score and the
  reverse flow animated. The purpose was calibration: seeing
  \eqref{eq:pt-convolution} act on a multimodal density, and the reverse
  drift reassemble it.
  \item \emph{TM2: two frames, additive dynamics.} The first genuinely
  dynamic object: $x_1 = x_0 + c + \eta$ with a trimodal prior on $x_0$,
  constant drift $c$, and Gaussian innovation $\eta$. The joint density
  $p(x_0, x_1)$ is a product of the prior and the transition kernel, and
  the two-dimensional joint score is computable in closed form.
  Figure~\ref{fig:toymodel} shows the object this toy model was built
  around: the joint density and its score field, whose diagonal
  orientation is the visible fingerprint of the frame coupling.
  \item \emph{TM3: rotation observed through OU noise.} The
  deterministic-dynamics variant of the board models: a planar rotation
  per frame, with the diffusion channel acting on the coordinates.
  \item \emph{TM4: bounded deterministic dynamics.} A bouncing cube:
  a deterministic bounded trajectory, chosen to probe score fields
  for dynamics with reflections, where the stationary law is far from
  Gaussian.
  \item \emph{TM5: two clocks.} Two coupled periodic phases advancing at
  different rates, the minimal model of a system with two internal
  frequencies.
  \item \emph{TM6: the circle prior.} The static circle distribution
  itself, isolating the geometry of the score of a curved data manifold
  from any internal dynamics.
  \item \emph{TM7: the two-frame joint score laboratory.} The synthesis
  of TM1--TM6 built after the correction described next: an interactive
  computation of the \emph{joint} two-frame score, with the coupling
  between frames on explicit display.
\end{itemize}

\begin{figure}[t]\centering
\includegraphics[width=\linewidth]{fig_toymodel_score.pdf}
\caption{The two-frame toy model (TM2/TM7): $x_1 = x_0 + c + \eta$ with a
trimodal Gaussian-mixture prior on $x_0$, drift $c = 1$, innovation
standard deviation $\sigma = 0.45$, after diffusion time $t = 0.15$.
Left: the noised joint density $p_t(x_0, x_1)$. Right: the joint score
field $\nabla \log p_t$ (arrows) over the same density (contours). Both
panels are computed from the closed-form mixture expressions, on a grid;
the arrows point towards the modes and tilt along the diagonal, the
visible effect of the coupling between the two frames.}
\label{fig:toymodel}
\end{figure}

The toy models delivered the geometric intuition they were built for,
and one lasting technical observation: everywhere a closed form
existed, the score's cross-frame components were controlled by the
\emph{conditional} structure of the joint law, not by the marginal law
of single frames. The toy models also exposed a conceptual error in
how the project was computing scores.

\section{The joint-score correction}\label{sec:dev-correction}

In the early experiments, scores were computed \emph{per frame}: each
frame's noisy marginal was treated as a one-dimensional distribution,
its score computed independently, and the trajectory treated in effect
as $K$ separate samples. Supervision corrected the object of study: a
sequence is \emph{one} sample from a joint law, and the quantity that
drives the reverse dynamics of the sequence is the joint score of the
whole noisy trajectory, i.e.\ the gradient of $\log P_t(x_0, \dots,
x_{K-1})$, whose $k$-th component is a conditional expectation over the
\emph{full} posterior of the clean frames, as in
\eqref{eq:tweedie-joint}.

For the stationary models of this thesis the correction changes the
object of study completely. Under the variance-preserving channel every
noisy frame is marginally $\Nd(0, 1)$ at every diffusion time, so the
per-frame marginal score is $-x_k$: independent of $t$, independent of
the dynamics, identical for every model in this thesis. All information
about the temporal structure lives in the joint object.
Chapter~\ref{ch:gaussian} exhibits the contrast in closed form at
$K = 2$. The correction also reframed the propagator hypothesis
\eqref{eq:dev-propagator}: the structure worth studying is not a map
between per-frame scores but the coupling pattern of the joint score,
which is a statement about a matrix (in the Gaussian case) or about a
posterior (in general).

\section{The minimal model}\label{sec:dev-minimal}

The joint-score reframing selected the model of the research chapters.
Among the candidate dynamic objects (circles, rotations, clocks), the
stationary Gaussian AR(1) chain under coordinatewise OU noise is the
minimal one in which the joint score is nontrivial yet exactly
computable: one coupling parameter $\alpha$ for the internal dynamics
against one noise parameter $t$ for the corruption, linear algebra
available at every step, and every structural question (coupling range,
locality, message passing) posable in closed form. The nonlinear board
models were set aside at that point, deliberately and explicitly, as targets
for after the solvable case was understood. Chapter~\ref{ch:ring} takes up the
first of them.

